\documentclass[journal]{IEEEtran}
\usepackage{booktabs}
\usepackage{tabularx}
\usepackage{array}
\usepackage[hidelinks]{hyperref}
\title{Supplementary File S6: Integrated Simulator Run}
\author{}
\begin{document}
\maketitle

\noindent This supplement reports the run described in Section~V. The
accompanying ZIP contains selected machine-readable records, source tables,
and figures; it is not a runnable system image.

\section{Run Conditions and Phase Checks}
The run used a deterministic cast-film simulator and tool adapter. No
language-model loop or physical equipment was used. The platform and
simulator were already running, and the run reused a shared data directory.
Its identifier was \texttt{20261031034253-1dbo}; the command was
\texttt{node bench/pipeline.mjs --profile integrated --seed 42 --cl-seeds 3}.
The recorded runtime was Node.js 24.19.0 on Windows x64, and the source
commit was \texttt{60fa983e}. All measurements refer to that snapshot; they
do not validate the repository's current HEAD.

All 17 phase gates passed. Of the 77 underlying checks, 76 passed, one raised
a warning, and none failed. The warning came from the BOPET semantic-card
check in P10; the phase gate itself passed. Table~\ref{tab:s6-stages} groups
the checks by pipeline stage.

\begin{table}[!t]
\centering
\caption{Grouped check outcomes from the integrated run.}
\label{tab:s6-stages}
\footnotesize
\setlength{\tabcolsep}{3pt}
\begin{tabularx}{\columnwidth}{@{}>{\raggedright\arraybackslash}p{2.65cm}r r r X@{}}
\toprule
Pipeline stages & Pass & Warn & Fail & Step purpose\\
\midrule
P0--P2 & 11 & 0 & 0 & Bootstrap, model reference, and provisioning\\
P3, P4, P4m, P4b--P4f & 30 & 0 & 0 & Acquisition, operation, governance, and records\\
P6--P7 & 7 & 0 & 0 & Seeded objective and multimodal acquisition checks\\
P8, P8b & 16 & 0 & 0 & Second-scene check and bounded-backstop drill\\
P9 & 3 & 0 & 0 & Platform subsystem checks\\
P10 & 9 & 1 & 0 & BOPET line provisioning and scripted task path\\
\midrule
Total & 76 & 1 & 0 & 17 phase gates passed\\
\bottomrule
\end{tabularx}
\end{table}

\section{Acquisition, Write Probes, and Latency}
The primary cast-film scene used five simulated endpoints: Modbus TCP,
OPC UA, MQTT, HTTP, and Modbus RTU. Five acquisition nodes returned 54 samples.
Modbus RTU was an acquisition
endpoint without a write-control node. Each of the other four protocol paths
received six negative probes, all of which the configured write-admission
checks rejected. Each writable path also recorded one accepted legal write.
The run did not record all legal-write attempts, so it does not support an
estimate of the false-block rate.

The \texttt{film-line} configuration was checked in the same run across five
endpoints: Modbus TCP, OPC UA, MQTT, Modbus RTU, and HTTP. They returned 12,
12, 5, 11, and 11 samples, respectively (51 total). Each of its three
writable paths recorded one accepted legal write and rejected three negative
probes (9/9 total). The report recorded no false block or code change. This
was an in-run configuration check, not an independent deployment or
repetition, and it does not estimate a general false-block rate.

\begin{table}[!t]
\centering
\caption{Endpoint checks for the in-run \texttt{film-line} configuration.}
\label{tab:s6-second-scene}
\footnotesize
\setlength{\tabcolsep}{3pt}
\begin{tabularx}{\columnwidth}{@{}>{\raggedright\arraybackslash}p{1.45cm}r r X@{}}
\toprule
Protocol & Samples & Rejected probes & Legal-write check\\
\midrule
Modbus TCP & 12 & 3/3 & One accepted\\
OPC UA & 12 & 3/3 & One accepted\\
MQTT & 5 & 3/3 & One accepted\\
Modbus RTU & 11 & -- & Acquisition only\\
HTTP & 11 & -- & Acquisition only\\
\midrule
Total & 51 & 9/9 & Three accepted checks\\
\bottomrule
\end{tabularx}
\end{table}

\begin{table}[!t]
\centering
\caption{Protocol-level measurements recorded in the integrated run.}
\label{tab:s6-protocol}
\footnotesize
\setlength{\tabcolsep}{3pt}
\begin{tabularx}{\columnwidth}{@{}>{\raggedright\arraybackslash}p{1.35cm}r r X@{}}
\toprule
Protocol & Samples & Rejected probes & Reported $p_{50}/p_{95}$ (ms)\\
\midrule
Modbus TCP & 9 & 6/6 & 47.673/48.348\\
OPC UA & 11 & 6/6 & 21.494/24.175\\
MQTT & 6 & 6/6 & 23.015/94.108\\
HTTP & 13 & 6/6 & 34.414/47.784\\
Modbus RTU & 15 & -- & --\\
\bottomrule
\end{tabularx}
\end{table}

The run stored these latency summaries by protocol. Its log mentions
approximately six writes but does not retain the exact per-protocol
denominators or event-level timings. The summary's 31.649\,ms value is the
unweighted mean of the four protocol-level $p_{50}$ values, not a pooled
median. These values are descriptive; they do not support adapter rankings or
population-level tail estimates.

The same run exercised one approval-enabled write path (P4c). A Line 2
request for configured value 145.5 remained pending until approval returned
status 200. The pending write was then released; the report records a simulated
PLC readback of 146 and a 452-ms pending-to-approval interval, including
polling. No approval-disabled condition was run, so this check documents the
enabled path but does not estimate the effect of approval.

\section{Seeded Objective and BOPET Task}
The three fixed-seed cases used the simulator's offline reference
$J^*=89.894$. The simulator score is
$J=55J_{\mathrm{thick}}+25J_{\mathrm{quality}}+8J_{\mathrm{energy}}+7J_{\mathrm{throughput}}$.
The thickness term scores deviation from $50\pm2\,\mu\mathrm{m}$ with soft
degradation outside the band; the quality term reflects defect rate, the
energy term screw speed, and the throughput term line speed. The score is
defined only for melt temperature in $195$--$225\,^{\circ}\mathrm{C}$ and
pressure no greater than $22\,\mathrm{MPa}$. The runner defines $J_{end}$ as
the mean of the last two valid trajectory scores, or as the sole score when
the stopping rule ends a case before a second score is available.
Table~\ref{tab:s6-seeds} also lists $J_{last}$, the final individual score.
Seed 43 began inside the configured stopping region and ended with no
iterations or writes. It is a zero-update case and provides no evidence of
iterative convergence. Seeds 42 and 44 each used two writes.

\begin{table}[!t]
\centering
\caption{Fixed-seed simulator objective values. $J_0$ is the initial score;
$J_{end}$ is the mean of the last two valid trajectory scores, or the sole valid score if the seed stops
before a second score is available; $J_{last}$ is the final per-iteration
score retained in the trajectory.}
\label{tab:s6-seeds}
\footnotesize
\setlength{\tabcolsep}{3pt}
\begin{tabular}{@{}r r r r r r r r@{}}
\toprule
Seed & $J_0$ & $J_{end}$ & $J_{last}$ & $J^*$ & Ratio & Iter. & Writes\\
\midrule
42 & 66.271 & 88.775 & 88.853 & 89.894 & 98.8\% & 2 & 2\\
43 & 86.542 & 86.542 & 86.542 & 89.894 & 96.3\% & 0 & 0\\
44 & 63.981 & 88.384 & 88.540 & 89.894 & 98.3\% & 2 & 2\\
\bottomrule
\end{tabular}
\end{table}

The BOPET thickness readings were 28.05, 27.98, 27.98, and
$24.84\,\mu\mathrm{m}$, from the initial state through the three writes. The
final reading was within the configured $25.0\pm0.7\,\mu\mathrm{m}$ band.
The task board marked the task COMPLETED, but its associated
parameter-change/optimization record remained open and unjudged. This is an
integration observation; it does not establish a completed optimization or
show that the writes caused the measured change.

\section{Provenance and Interpretation}
The source run did not include a contemporaneous SHA-256 manifest. The ZIP
hashes verify the delivered files, not the runtime state at the time of
execution.

\end{document}
